\section{Introduction}

The binary distinction between ``existence'' and ``non-existence'' appears empirically as a two-valued discrimination, yet can theoretically be viewed as the result of projecting a single ontological object via a readout protocol. If observation is modeled as orthogonal cuts with finite resolution, the observed binary quantity need neither be presupposed as an ontologically primitive entity nor require appeal to external assertions beyond mathematical structure; its appearance can be generated through three operational steps: ``scan--project--encode'', and filtered by stability requirements to select repeatable readout structures.

The goal of this paper is to provide a model that simultaneously satisfies the following conditions:
\begin{enumerate}
  \item \textbf{Computability:} All finite-resolution objects can be exhaustively enumerated or generated by explicit algorithms on finite candidate classes;
  \item \textbf{Auditability:} Definitions, theorems, and interface assumptions are strictly stratified to avoid semantic jumps in proof chains;
  \item \textbf{Falsifiability:} Model outputs appear as discrete verifiable invariants (counts, spectra, degeneracies, boundary structures, statistical test protocols).
\end{enumerate}

Under this framework, the special character of the golden branch is not stated as metaphysical superiority but rather characterized through uniformity certificates under finite resources and minimal syntactic complexity: Kronecker orbits generated by slopes related to the golden ratio exhibit optimal (or near-optimal) low-discrepancy constants under finite sampling; the corresponding binary mechanistic words belong to the Fibonacci/Sturmian class with minimal factor complexity; syntactic filtering for window legality yields stable type spaces with Fibonacci-like growth. Thus, ``golden existence state'' can be defined as a normalized state on the stable subspace, and ``existence/non-existence'' emerges as the binary result of projection readout.
